\begin{abstract}
Process verifiers are often evaluated by ordinary valid-vs.-invalid
discrimination. In answer-visible settings, however, those metrics confound
local process grounding with answer-consistency shortcuts. We introduce
answer-matched counterfactual audits for answer-visible process verifiers,
centered on answer-matched counterfactual discrimination (\amcd) and
answer-swap sensitivity (\ass), and instantiate them on executable
matched-quartet benchmarks in algebra, graph path, and blocksworld. Across
prompt judges, frozen-head baselines, diagnostic repair objectives, and
stronger rerankers, ordinary metrics materially overstate the reliability of
answer-visible verification. Judge-style evaluators are strongly
answer-sensitive, and the mechanism analyses are consistent with lower local
\amcd{} accompanying worse downstream utility. On the deployment benchmark
\textsc{Craft-Deploy}, a masked \texttt{Qwen3-Reranker-8B} is the strongest
model, outperforming visible reranking and frozen-head baselines on \amcd{},
exploitability, robustness, and seed-averaged reproduction. On the cleaner
audit benchmark \textsc{Craft-Clean}, masking still improves
process-faithfulness while utility no longer separates, indicating a cleaner
faithfulness benchmark rather than a direct deployment surrogate. Minimal
repairs act as diagnostic interventions: they move the
faithfulness--utility frontier but do not dominate strong masked baselines.
On the external-source \textsc{ProcessBench} package, the same conservative
read also holds: a step-only frozen baseline is currently strongest,
answer-only swaps still move visible scores far more than masked scores, and a
fully constructed local-repair subset still does not restore a visible
advantage.
The operational conclusion is simple: answer-visible process verifiers should
be audited with answer-matched counterfactuals, and the safest current
deployment choice is to remove direct answer access.
\end{abstract}
